\section{The canonical proxy cost}\label{app:objective}

This appendix states the score exactly as the pinned evaluator computes it. The
formulation follows Circuit Training \citep{mirhoseini2021graph}; the evaluator is the
open TILOS \texttt{PlacementCost} class in \path{plc_client_os.py}
(MacroPlacement revision \texttt{45a721d}); the challenge harness wraps it in
\path{compute_proxy_cost} (\path{objective.py}). Function names below refer
to those two files. Every rule was checked against that source for this report.

\paragraph{Notation.} The canvas is $[0,L_x]\times[0,L_y]$, divided into a grid of
$n_c$ columns and $n_r$ rows with $G=n_cn_r$ cells of width $\Delta_x=L_x/n_c$ and
height $\Delta_y=L_y/n_r$. Object $i$ (hard macro or soft cluster) has center
$p_i=(x_i,y_i)$, width $w_i$ and height $h_i$. A hard-macro pin sits at its
macro's center plus a fixed offset; a soft-cluster pin sits at the cluster center;
a port (an I/O terminal on the canvas edge) has a fixed position. The grid
dimensions, the canvas size and the routing constants are read from the header of
each design's \path{initial.plc} (\path{restore_placement} with comment
parsing). All 17 IBM headers carry the same constants: horizontal and vertical
routing capacities $r_H=65.957$ and $r_V=106.957$ routes per $\mu$m, macro routing
allocations $a_H=30.304$ and $a_V=71.304$ per $\mu$m, and smoothing range $s=2$
cells. The header's overlap threshold (0.004) does not enter the score.

\begin{table}[htbp]
\centering
\footnotesize
\caption{Benchmark statistics of the 17 IBM designs, parsed from the challenge inputs with the canonical evaluator's loader. Nets is the evaluator's net count; the grid is the density and congestion grid.}
\label{tab:benchmarks}
\begin{tabular}{lrrrcc}
\toprule
Design & \shortstack{Movable\\hard macros} & \shortstack{Soft\\clusters} & Nets & \shortstack{Canvas $W \times H$\\($\mu$m)} & \shortstack{Grid\\columns $\times$ rows} \\
\midrule
IBM01 & 246 & 894 & 7{,}269 & 22.95 $\times$ 23.04 & 45 $\times$ 41 \\
IBM02 & 271 & 1{,}075 & 10{,}944 & 32.56 $\times$ 32.48 & 30 $\times$ 27 \\
IBM03 & 290 & 1{,}148 & 10{,}247 & 35.12 $\times$ 35.04 & 32 $\times$ 29 \\
IBM04 & 295 & 1{,}085 & 13{,}555 & 34.09 $\times$ 34.08 & 31 $\times$ 30 \\
IBM06 & 178 & 900 & 14{,}085 & 32.65 $\times$ 32.64 & 31 $\times$ 28 \\
IBM07 & 291 & 1{,}040 & 19{,}920 & 38.51 $\times$ 38.40 & 35 $\times$ 32 \\
IBM08 & 301 & 1{,}030 & 20{,}694 & 41.05 $\times$ 40.96 & 38 $\times$ 34 \\
IBM09 & 253 & 1{,}048 & 18{,}341 & 46.74 $\times$ 46.88 & 36 $\times$ 38 \\
IBM10 & 786 & 1{,}982 & 39{,}140 & 77.05 $\times$ 77.12 & 55 $\times$ 41 \\
IBM11 & 373 & 1{,}195 & 25{,}273 & 51.53 $\times$ 51.52 & 39 $\times$ 45 \\
IBM12 & 651 & 1{,}985 & 40{,}996 & 67.97 $\times$ 68.00 & 47 $\times$ 47 \\
IBM13 & 424 & 1{,}301 & 28{,}201 & 55.94 $\times$ 56.00 & 43 $\times$ 43 \\
IBM14 & 614 & 1{,}529 & 57{,}333 & 59.85 $\times$ 60.00 & 49 $\times$ 44 \\
IBM15 & 393 & 1{,}138 & 46{,}467 & 67.64 $\times$ 67.52 & 57 $\times$ 38 \\
IBM16 & 458 & 1{,}315 & 67{,}118 & 81.03 $\times$ 81.12 & 45 $\times$ 48 \\
IBM17 & 760 & 1{,}844 & 83{,}132 & 72.60 $\times$ 72.64 & 51 $\times$ 44 \\
IBM18 & 285 & 1{,}029 & 53{,}763 & 64.82 $\times$ 64.96 & 55 $\times$ 39 \\
\bottomrule
\end{tabular}
\par\smallskip
\begin{minipage}{0.92\linewidth}\footnotesize No design has fixed hard macros. IBM05 is excluded because it has no macros; the benchmark set used by the challenge omits it. The macro and net counts summarized in the challenge README agree with these parsed values only for IBM01.\end{minipage}
\end{table}


\paragraph{Weights.} The score is
\begin{equation}
  J(p)=W(p)+\tfrac12\,D(p)+\tfrac12\,C(p),
  \label{eq:app-proxy}
\end{equation}
with the three components defined below (\path{compute_proxy_cost}). The density
component already contains a factor $\tfrac12$, so the raw density tail enters $J$
with coefficient $\tfrac14$.

\paragraph{Wirelength.} Each pin or port with a non-empty list of sinks is the
driver of one net $n$, whose pins are the driver and its sinks. The net weight
$\omega_n$ is the driver's weight attribute (default~1; ports count with weight~1).
With the half-perimeter wirelength
\[
 \mathrm{HPWL}_n=\bigl(\max_{q\in n}x_q-\min_{q\in n}x_q\bigr)
 +\bigl(\max_{q\in n}y_q-\min_{q\in n}y_q\bigr),
\]
the component is
\begin{equation}
  W(p)=\frac{\sum_n\omega_n\,\mathrm{HPWL}_n(p)}{(L_x+L_y)\,N_\omega},
  \qquad N_\omega=\sum_n\omega_n
\end{equation}
(\path{get_cost} and \path{get_wirelength}). The normalizer $N_\omega$ is the
evaluator's weighted net counter \path{net_cnt}, accumulated while the netlist is
parsed; it is not the number of entries in the net dictionary.

\paragraph{Density.} Let $R_i$ be the rectangle of object $i$ and $B_b$ cell $b$.
Hard macros and soft clusters both contribute, through exact intersection areas:
\begin{equation}
  \rho_b(p)=\frac{1}{\Delta_x\Delta_y}\sum_i\bigl|R_i\cap B_b\bigr|
\end{equation}
(\path{get_grid_cells_density}). For $G\ge10$, let $k=\lfloor0.1\,G\rfloor$.
The evaluator sorts the nonzero loads in decreasing order, sums the largest $k$ of
them (or all of them, if fewer than $k$ are nonzero) and divides by $k$:
\begin{equation}
  D(p)=\frac12\cdot\frac1k\sum_{j=1}^{\min(k,n_{\mathrm{occ}})}\rho_{(j)},
\end{equation}
where $n_{\mathrm{occ}}$ counts nonzero cells. Dropping zeros while keeping the
denominator $k$ is the same as padding with zero loads. For a grid with $G<10$
cells the rule changes: $D$ is half the mean over the occupied cells
only (\path{get_density_cost}). Every IBM grid has between 810 and 2,255 cells,
so the small-grid rule never applies to the challenge designs.

\paragraph{Congestion: net routing.} Congestion uses a pattern-routing estimate on
the same grid (\path{get_routing}). Each pin is mapped to the cell containing it,
$(\lfloor y/\Delta_y\rfloor,\lfloor x/\Delta_x\rfloor)$; the challenge harness patches
this lookup to clamp indices into the grid, so off-canvas pins land in border cells.
For net $n$, let $S_n$ be the set of distinct cells of its pins and $s_n$ the
driver's cell. Each cell along a route segment gains demand $\bar\omega_n$, equal to
the driver weight when it exceeds~1 and to~1 otherwise:
\begin{itemize}
  \item $|S_n|=1$: no demand.
  \item $|S_n|=2$: an L-shaped route. Horizontal demand runs along the driver's row
        from the smaller to the larger column, excluding the last; vertical demand
        runs along the sink's column from the smaller to the larger row, excluding
        the last (\path{__two_pin_net_routing}).
  \item $|S_n|=3$: an L- or T-shaped Steiner route chosen from the relative order of
        the three cells (\path{__three_pin_net_routing}).
  \item $|S_n|>3$: a star of two-pin routes from $s_n$ to every other cell of
        $S_n$ (\path{__split_net}).
\end{itemize}
The accumulated vertical and horizontal counts are divided by the per-cell
capacities $c_V=\Delta_x r_V$ and $c_H=\Delta_y r_H$.

\paragraph{Congestion: smoothing and macro blockage.} The normalized net demand is
smoothed along one axis only (\path{__smooth_routing_cong}). Each vertical
value is spread uniformly over the cells within $\pm s$ columns of the same row, and
each horizontal value over the cells within $\pm s$ rows of the same column; at the
grid boundary the window is clipped and the value is divided by the clipped window
size, so smoothing conserves demand. Hard macros then add routing blockage
(\path{__macro_route_over_grid_cell}). For each cell a hard macro covers,
the vertical map gains (horizontal overlap length)$\times a_V$ and the horizontal
map gains (vertical overlap length)$\times a_H$, both divided by the same capacities.
A macro that fully covers a cell therefore adds $a_V/r_V=0.667$ to its vertical
value and $a_H/r_H=0.459$ to its horizontal value. One edge correction applies: if a
macro spans several rows and covers its lowest or highest row only partly, the
vertical blockage of its entire top row is removed; if it spans several columns and
covers its first or last column only partly, the horizontal blockage of its entire
rightmost column is removed. The blockage is added after smoothing, so it is not
smoothed:
\begin{equation}
  V=\mathcal S_{\mathrm{col}}(V_{\mathrm{net}})+V_{\mathrm{mac}},\qquad
  H=\mathcal S_{\mathrm{row}}(H_{\mathrm{net}})+H_{\mathrm{mac}}.
\end{equation}

\paragraph{Congestion: tail statistic.} The score treats the two maps as one list of
$2G$ values, all vertical values followed by all horizontal ones, and takes the mean of
its largest $\lfloor0.05\cdot2G\rfloor$ entries, zeros included (\texttt{abu} with fraction
0.05, called by \path{get_congestion_cost}). If that count is zero, the maximum
is used. The evaluator's separate per-direction functions
(\path{get_V_congestion_cost}, \path{get_H_congestion_cost}) do not enter
the score.

\paragraph{What the definition implies for search.} All three components are
piecewise: HPWL through pin extrema, density through rectangle--cell
intersections, congestion through cell indices, routing patterns and the two tail
selections. A move of one object changes the pins, nets and cells it touches, but
both tails are recomputed over the whole grid, so the cost of scoring a move remains
linear in $G$ (Remark~\ref{rem:move-cost}). Congestion depends on cell indices
rather than exact coordinates; a small move that keeps every pin in its cell leaves
the net-routing demand unchanged, and the score jumps when a pin crosses a grid line,
so the routing part is piecewise constant rather than merely kinked.


\paragraph{The benchmarks.} The 17 designs are the ICCAD04 IBM mixed-size
benchmarks \citep{adya2004unification} as converted by the TILOS MacroPlacement
project \citep{macroplacement}; Table~\ref{tab:benchmarks} lists their statistics,
parsed from the pinned inputs. The conversion groups each design's standard cells
into \emph{soft clusters} with the Circuit Training grouping flow
\citep{mirhoseini2021graph}, starting from a RePlAce initial placement
\citep{cheng2019replace}. A soft cluster is a fixed-size rectangle that may
overlap other clusters and macros; it carries pins and so contributes to all three
components. The score therefore depends on where the soft clusters sit, and
Thinkorplace moves them together with the hard macros. They exist only inside the
proxy: the organizers lock their sizes and state that the soft-macro abstraction
does not carry over to place-and-route. No IBM design has fixed hard macros, and
hard macros and soft clusters together occupy 80\% of the canvas area in every
design. The suite skips IBM05 because, as the TILOS ICCAD04 testcase notes state,
that design has no macros.
